\newpage
\section{问题重述}

\subsection{问题背景}
输电线路是电力系统的"动脉"，其安全运行关乎电网稳定。传统的人工巡检效率低、风险高、覆盖不全。无人机巡检具有灵活、高效、成本低的优点，已成为输电线路巡检的主流手段。但无人机巡检路径规划涉及复杂的约束条件和优化目标，是典型的组合优化问题。

\subsection{问题提出}
\textbf{问题一}：将输电线路网络抽象为图结构，建立图论模型描述航线网络，分析网络的拓扑特性。

\textbf{问题二}：构建无人机巡检路径优化模型，以巡检效率最高、能耗最低为目标，考虑电池容量、飞行高度、拍摄角度等约束。

\textbf{问题三}：设计改进的NSGA-II算法求解多目标优化问题，并分析算法参数对求解效果的影响。

\subsection{研究意义}
为无人机智能巡检提供路径规划方法，提升输电线路运维效率，降低人工巡检风险。
